\emph{By Chris Harman}

Few questions have produced more bitterness in Marxist circles than that of the relation between the party and the class. More heat has probably been generated in acrimonious disputes over this subject than any other. In generation after generation the same epithets are thrown about---``bureaucrat'', ``substitutionist'', ``elitist'', ``autocrat''.

Yet the principles underlying such debate have usually been confused. This despite the importance of the issues involved. For instance, the split between Bolsheviks and Mensheviks that occurred over the nature of the organisation of the party in 1903 found many of those who were to be on the opposite side of the barricades to Lenin in 1917 in his faction (for instance, Plekhanov), while against him were revolutionaries of the stature of Trotsky and Rosa Luxemburg.

Nor was this confusion an isolated incident. It has been a continuous feature of revolutionary discussion. It is worth recalling Trotsky's remarks, at the second Congress of the Comintern, in reply to Paul Levi's contention that the mass of workers of Europe and America understood the need for a party. Trotsky points out that the situation is much more complex than this:

\blockquote{If the question is posed in the abstract then I see Scheidemann on the one side and, on the other, American or French or Spanish syndicalists who not only wish to fight against the bourgeoisie, but who, unlike Scheidemann, really want to tear its head off---for this reason I say that I prefer to discuss with these Spanish, American or French comrades in order to prove to them that the party is indispensable for the fulfilment of the historical mission which is placed upon them\ldots I will try to prove this to them in a comradely way, on the basis of my own experience, and not by counterposing to them Scheidemann's long years of experience saying that for the majority the question has already been settled\ldots What is there in common between me and a Renaudel who excellently understands the need of the party, or an Albert Thomas and other gentlemen whom I do not even want to call ``comrades'' so as not to violate the rules of decency?\footnote{Leon Trotsky, \emph{The First Five Years of the Communist International}, Vol.~1, New York 1977, p.~98.}}

The difficulty to which Trotsky refers---that both Social Democrats and Bolsheviks refer to the ``need for a party'', although what they mean by this are quite distinct things---has been aggravated in the years since by the rise of Stalinism. The vocabulary of Bolshevism was taken over and used for purposes quite opposed to those who formulated it. Yet too often those who have continued in the revolutionary tradition opposed to both Stalinism and Social Democracy have not taken Trotsky's points in 1920 seriously. They have often relied on ``experience'' to prove the need for a party, although the experience is that of Stalinism and Social Democracy. Most of the discussion even in revolutionary circles is, as a consequence, discussion for or against basically Stalinist or Social-Democratic conceptions of organisation. The sort of organisational views developed implicitly in the writings and actions of Lenin are radically different to both these conceptions. This has been obscured by the Stalinist debasement of the theory and practice of the October revolution and the fact that the development of the Bolshevik Party took place under conditions of illegality and was often argued for in the language of orthodox Social Democracy.

\paragraph{The Social-Democratic View of the Relation of Party and Class}

The classical theories of Social Democracy---which were not fundamentally challenged by any of the Marxists before 1914---of necessity gave the party a central role in the development towards socialism. For this development was seen essentially as being through a continuous and smooth growth of working-class organisation and consciousness under capitalism. Even those Marxists, such as Kautsky, who rejected the idea that there could be a gradual transition to socialism accepted that what was needed for the present was continually to extend organisational strength and electoral following. The growth of the party was essential so as to ensure that when the transition to socialism inevitably came, whether through elections or through defensive violence by the working class, the party capable of taking over and forming the basis of the new state would exist.

The development of a mass working-class party is seen as being an inevitable corollary of the tendencies of capitalist development. ``Forever greater grows the number of proletarians, more gigantic the army of superfluous labourers, and sharper the opposition between exploiters and exploited'',\footnote{Karl Kautsky, \emph{The Erfurt Program}, Chicago 1910, p.~8.} crises ``naturally occur on an increasing scale'',\footnote{Ibid.} ``the majority of people sink ever deeper into want and misery'',\footnote{Ibid., p.~43.} ``the intervals of prosperity become ever shorter; the length of the crises ever longer''.\footnote{Ibid., p.~85.} This drives greater numbers of workers ``into instinctive opposition to the existing order''.\footnote{Ibid., p.~198.} Social Democracy, basing itself upon ``independent scientific investigation by bourgeois thinkers'',\footnote{Ibid., p.~198.} exists to raise the workers to the level where they have a ``clear insight into social laws''.\footnote{Ibid., p.~198.} Such a movement ``springing out of class antagonisms\ldots cannot meet with anything more than temporary defeats, and must ultimately win''.\footnote{Karl Kautsky, \emph{The Road to Power}, Chicago 1910, p.~24.} ``Revolutions are not made at will\ldots They come with inevitable necessity.'' The central mechanism involved in this development is that of parliamentary elections. ``We have no reason to believe that armed insurrection\ldots will play a central role nowadays.''\footnote{Karl Kautsky, \emph{Social Revolution}, p.~47.} Rather, ``it (parliament) is the most powerful lever that can be used to raise the proletariat out of its economic, social and moral degradation.''\footnote{Karl Kautsky, \emph{The Erfurt Program}, p.~188.} The uses of this by the working class makes ``parliamentarianism begin to change its character. It ceases to be a mere tool in the hands of the bourgeoisie.''\footnote{Ibid., p.~188.} In the long run such activities must lead to the organisation of the working class and to a situation where the socialist party has the majority and will form the government. ``\ldots (The Labour Party) must have for its purpose the conquest of the government in the interests of the class it represents. Economic development will lead naturally to the accomplishment of this purpose.''\footnote{Ibid., p.~189.}

What is central for the Social Democrat is that the party \emph{represents} the class. Outside of the party the worker has no consciousness. Indeed, Kautsky himself seemed to have an almost pathological fear of what the workers would do without the party and of the associated dangers of a ``premature'' revolution. Thus it had to be the party that takes power. Other forms of working-class organisation and activity can help, but must be subordinated to the bearer of political consciousness. ``This `direct action' of the unions can operate effectively only as an auxiliary and reinforcement to and not as a substitute for parliamentary action.''\footnote{Karl Kautsky, \emph{The Road to Power}, p.~95.}

\paragraph{The Revolutionary Left and Social-Democratic Theories}

No sense can be made of any of the discussions that took place in relation to questions of organisation of the party prior to 1917 without understanding that this Social-Democratic view of the relation of party and class was \emph{nowhere} explicitly challenged (except among the anarchists who rejected any notion of a party). Its assumptions were shared even by those, such as Rosa Luxemburg, who opposed orthodox Social Democracy from the point of view of mass working-class self-activity.

This was not a merely theoretical failing. It followed from the historical situation. The Paris Commune was the only experience then of working-class power, and that had been for a mere two months in a predominantly petty-bourgeois city. Even the 1905 revolution gave only the most embryonic expression of how a workers' state would in fact be organised. The fundamental forms of workers' power---the soviets, the workers' councils---were not recognised. Thus Trotsky, who had been President of the Petrograd Soviet in 1905, does not mention them in his analysis of the lessons of 1905, \emph{Results and Prospects}. Virtually alone in foreseeing the socialist content of the Russian revolution, Trotsky did not begin to see the form this would take. Not until the February revolution did the soviet become central in Lenin's writings and thoughts.

The revolutionary Left never fully accepted Kautsky's position of seeing the party as the direct forerunner of the workers' state. Luxemburg's writings, for instance, recognise the conservatism of the party and the need for the masses to go beyond and outside it from a very early stage.\footnote{Cf. both \emph{Organisational Questions of the Russian Social-Democracy} (published by her epigones under the title \emph{Leninism or Marxism}), and \emph{The Mass Strike, the Political Party and the Trade Unions}.} But there is never an explicit rejection of the official Social-Democratic position. Yet without the theoretical clarification of the relationship between the party and the class there could be no possibility of clarity over the question of the necessary internal organisation of the party. Without a rejection of the Social-Democratic model, there could not be the beginnings of a real discussion about revolutionary organisation.

This is most clearly the case with Rosa Luxemburg. It would be wrong to fall into the trap (carefully laid by both Stalinist and would-be followers of Luxemburg) of ascribing to her a theory of ``spontaneity'' that ignores the need for a party. Throughout her writings there is stress upon the need for a party and the positive role it must play:

\blockquote{In Russia, however, the Social-Democratic Party must make up by its own efforts an entire historical period. It must lead the Russian proletarians from their present ``atomised'' condition, which prolongs the autocratic regime, to a class organisation that would help them to become aware of their historical objectives and prepare them to struggle to achieve those objectives.\footnote{Rosa Luxemburg, \emph{Leninism or Marxism}, Ann Arbor 1962, p.~82.}}

\blockquote{The task of Social Democracy does not consist in the technical preparation and direction of mass strikes, but first and foremost in the political leadership of the whole movement.\footnote{Rosa Luxemburg, \emph{The Mass Strike}, p.~57.}}

\blockquote{The Social Democrats are the most enlightened, the most class-conscious vanguard of the proletariat. They cannot and dare not wait, in a fatalistic fashion with folded arms for the advent of the ``revolutionary situation''.\footnote{Ibid.}}

Yet there is a continual equivocation in Luxemburg's writings on the role of the party. She was concerned that the leading role of the party should not be too great---for she identified this as ``the prudent position of Social Democracy''.\footnote{Rosa Luxemburg, \emph{Leninism or Marxism}, p.~92.} She identified ``centralism'', which she saw as anyway necessary (``the Social Democracy is, as a rule, hostile to any manifestation of localism or federalism''\footnote{Ibid., p.~85.}) with the ``conservatism inherent in such an organ (i.e. the Central Committee)''.\footnote{Ibid., p.~94.} Such equivocation cannot be understood without taking account of the concrete situation Luxemburg was really concerned about. She was a leading member of the SPD, but always uneasy about its mode of operation. When she really wanted to illustrate the dangers of centralism it was to this that she referred:

\blockquote{The present tactical policy of the German Social Democracy has won universal esteem because it is supple as well as firm. This is a sign of the fine adaptation of our party to the conditions of a parliamentary regime\ldots However, the very perfection of this adaptation is already closing vaster horizons to our party.}

Brilliantly prophetic as this is of what was to happen in 1914, she does not begin to explain the origins of the increasing sclerosis and ritualism of the SPD, let alone indicate ways of fighting this. Conscious individuals and groups cannot resist this trend. For ``such inertia is due, to a large degree to the fact that it is inconvenient to define, within the vacuum of abstract hypotheses, the lines and forms of non-existent political situations''.\footnote{Ibid., p.~93.} Bureaucratisation of the party is seen as an inevitable phenomenon that only a limitation on the degree of cohesion and efficiency of the party can overcome. It is not a particular form of organisation and conscious direction, but organisation and conscious direction as such that limit the possibilities for the ``self-conscious movement of the majority in the interests of the majority''.

\blockquote{The unconscious comes before the conscious. The logic of history comes before the subjective logic of the human beings who participate in the historic process. The tendency is for the directing organs of the socialist party to play a conservative role.\footnote{Ibid., p.~93.}}

There is a correct and important element in this argument: the tendency for certain sorts of organisations to be unable to respond to a rapidly changing situation. One only has to think of the Maximalist wing of the Italian Socialist Party in 1919, the whole of the ``centre'' of the Second International in 1914, the Menshevik-Internationalists in 1917, or the KPD in 1923. Even the Bolshevik Party contained a very strong tendency to exhibit such conservatism. But Luxemburg, having made the diagnosis, makes no attempt to locate its source, except in epistemological generalities, or looks for organisational remedies. There is a strong fatalism in her hope that the ``unconscious'' will be able to correct the ``conscious''. Despite her superb sensitivity to the peculiar tempo of development of the mass movement---particularly in \emph{The Mass Strike}---she shies away from trying to work out a clear conception of the sort of political organisation that can harness such spontaneous developments. Paradoxically this most trenchant critic of bureaucratic ritualism and parliamentary cretinism argued in the 1903 debate for precisely that faction of the Russian party that was to be the most perfected historical embodiment of these failings: the Mensheviks. In Germany political opposition to Kautskyism, which already was developing at the turn of the century and was fully formed by 1910, did not take on concrete organisational forms for another five years.

Considerable parallels exist between Luxemburg's position and that which Trotsky adhered to up to 1917. He too is very aware of the danger of bureaucratic ritualism:

\blockquote{The work of agitation and organisation among the ranks of the proletariat has an internal inertia. The European Socialist Parties, particularly the largest of them, the German Social-Democratic Party, have developed an inertia in proportion as the great masses have embraced socialism and the more these masses have become organised and disciplined. As a consequence of this, Social Democracy as an organisation embodying the political experience of the proletariat may at a certain moment become a direct obstacle to open conflict between the workers and bourgeois reaction.\footnote{Leon Trotsky, \emph{Results and Prospects} (1906), in \emph{The Permanent Revolution and Results and Prospects}, London 1962, p.~246.}}

Again his revolutionary spirit leads him to distrust \emph{all} centralised organisation. Lenin's conception of the party can, according to Trotsky in 1904, only lead to the situation in which:

\blockquote{The organisation of the Party substitutes itself for the party as a whole; then the Central Committee substitutes itself for the organisation; and finally the ``dictator'' substitutes himself for the Central Committee.\footnote{Quoted in I. Deutscher, \emph{The Prophet Armed}, London 1954, pp.~92--93.}}

But for Trotsky the real problems of working-class power can only be solved

\blockquote{by way of systematic struggle between\ldots many trends inside socialism, trends which will inevitably emerge as soon as the proletarian dictatorship poses tens and hundreds of new\ldots problems. No strong ``domineering'' organisation will be able to suppress these trends and controversies\ldots\footnote{Ibid.}}

Yet Trotsky's fear of organisational rigidity leads him also to support that tendency in the inner-party struggle in Russia which was historically to prove itself most frightened by the spontaneity of mass action. Although he was to become increasingly alienated from the Mensheviks politically, he did not begin to build up an organisation in opposition to them until very late. Whether he was correct or not in his criticisms of Lenin in 1904 (and we believe he was wrong), he was only able to become an effective historical actor in 1917 by joining Lenin's party.

If organisation does produce bureaucracy and inertia, Luxemburg and the young Trotsky were undoubtedly right about the need to limit the aspirations towards centralism and cohesion among revolutionaries. But it is important to accept all the consequences of this position. The most important must be a historical fatalism. Individuals can struggle among the working class for their ideas, and these ideas can be important in giving workers the necessary consciousness and confidence to fight for their own liberation. But revolutionaries can never build the organisation capable of giving them effectiveness and cohesion in action comparable to that of those who implicitly accept present ideologies. For to do so is inevitably to limit the self-activity of the masses, the ``unconscious'' that precedes the ``conscious''. The result must be to wait for ``spontaneous'' developments among the masses. In the meantime one might as well put up with the organisations that exist at present, even if one disagrees with them politically, as being the best possible, as being the maximum present expression of the spontaneous development of the masses.

\paragraph{Lenin and Gramsci on the Party and the Class}

In the writings of Lenin there is an ever-present implicit recognition of the problems that worry Luxemburg and Trotsky so much. But there is not the same fatalistic succumbing to them. There is an increasing recognition that it is not organisation as such, but particular forms and aspects of organisation that give rise to these. Not until the First World War and then the events in 1917 gave an acute expression to the faults of old forms of organisation did Lenin begin to give explicit notice of the radically new conceptions he himself was developing. Even then these were not fully developed. The destruction of the Russian working class, the collapse of any meaningful soviet system (i.e. one based upon real workers' councils), and the rise of Stalinism, smothered the renovation of socialist theory. The bureaucracy that arose with the decimation and demoralisation of the working class took over the theoretical foundations of the revolution, to distort them into an ideology justifying its own interests and crimes. Lenin's view of what the party is and how it should function in relation to the class and its institutions was no sooner defined as against older Social-Democratic conceptions with any clarity than it was again obscured by a new Stalinist ideology.

Many of Lenin's conceptions are, however, taken up and given clear and coherent theoretical form by the Italian Antonio Gramsci.\footnote{Antonio Gramsci, \emph{The Prison Notebooks}.} Throughout Lenin's writings are two intertwined and complementary conceptions, which to the superficial observer seem contradictory. Firstly there is continual stress on the possibilities of sudden transformations of working-class consciousness, on the unexpected upsurge that characterises working-class self-activity, on deep-rooted instincts in the working class that lead it to begin to reject habits of deference and subservience:

\blockquote{In the history of revolutions there come to light contradictions that have ripened for decades and centuries. Life becomes unusually eventful. The masses, which have always stood in the shade and therefore have often been despised by superficial observers, enter the political arena as active combatants\ldots Nothing will ever compare in importance with this direct training that the masses and the classes receive in the course of the revolutionary struggle itself.\footnote{V. I. Lenin, ``Lecture on the 1905 Revolution.''}}

\blockquote{We are able to appreciate the importance of the slow, steady and often imperceptible work of political education which Social Democrats have always conducted and always will conduct. But we must not allow what in the present circumstances would be still more dangerous---a lack of faith in the powers of the people. We must remember what a tremendous educational and organisational power the revolution has, when mighty historical events force the man in the street out of his remote garret or basement corner, and make a citizen of him. Months of revolution sometimes educate citizens more quickly and fully than decades of political stagnation.\footnote{V. I. Lenin, ``Two Tactics of Social-Democracy in the Democratic Revolution.''}}

\blockquote{The working class is instinctively, spontaneously Social Democratic.\footnote{V. I. Lenin, \emph{What Is to Be Done?}}}

\blockquote{The special condition of the proletariat in capitalistic society leads to a striving of workers for socialism; a union of them with the Socialist Party bursts forth with a spontaneous force in the very early stages of the movement.\footnote{Ibid.}}

Even in the worst months after the outbreak of war in 1914 he could write:

\blockquote{The objective war-created situation\ldots is inevitably engendering revolutionary sentiments; it is tempering and enlightening all the finest and most class-conscious proletarians. A sudden change in the mood of the masses is not only possible, but is becoming more and more probable\ldots\footnote{V. I. Lenin, \emph{Collected Works}, Vol.~21, p.~80.}}

In 1917 this faith in the masses leads him in April and in August-September into conflict with his own party. As Trotsky wrote, ``Lenin said more than once that the masses are to the Left of the party. He knew the party was to the Left of its own upper layer of `old Bolsheviks'.''\footnote{Leon Trotsky, \emph{History of the Russian Revolution}.} In relation to the ``Democratic Conference'' he could write:

\blockquote{We must draw the masses into the discussion of this question. Class-conscious workers must take the matter into their own hands, organise the discussion and exert pressure on ``those at the top''.\footnote{V. I. Lenin, \emph{Collected Works}, Vol.~26.}}

There is, however, a second fundamental element in Lenin's thought and practice: the stress on the role of theory and of the party as the bearer of this. The most well known recognition of this occurs in \emph{What Is to Be Done?} when Lenin writes that ``Without revolutionary theory there can be no revolutionary practice.'' But it is the theme that recurs at every stage in his activities, not only in 1903, but also in 1905 and 1917 at exactly the same time that he was cursing the failure of the party to respond to the radicalisation of the masses. And for him the party is something very different from the mass organisations of the whole class. It is always a vanguard organisation, membership of which requires a dedication not to be found in most workers. (But this does not mean that Lenin ever wanted an organisation only of professional revolutionaries.) This might seem a clear contradiction, particularly as in 1903 Lenin uses arguments drawn from Kautsky which imply that only the party can imbue the class with a socialist consciousness, while later he refers to the class being more ``to the Left'' than the party. In fact, however, to see a contradiction here is to fail to understand the fundamentals of Lenin's thinking on these issues. For the real theoretical basis for his argument on the party is not that the working class is incapable on its own of coming to theoretical socialist consciousness. This he admits at the second congress of the RSDLP when he denies that ``Lenin takes no account whatever of the fact that the workers too have a share in the formation of an ideology'' and adds that ``\ldots The `economists' have gone to one extreme. To straighten matters out somebody had to pull in the other direction---and that is what I have done.''\footnote{V. I. Lenin, \emph{Collected Works}, Vol.~6.}

The real basis for his argument is that the level of consciousness in the working class is never uniform. However rapidly the mass of workers learn in a revolutionary situation, some sections will still be more advanced than others. To merely take delight in the spontaneous transformation is to accept uncritically whatever transitory products this throws up. But these reflect the backwardness of the class as well as its movement forward, its situation in bourgeois society as well as its potentiality for further development so as to make a revolution. Workers are not automatons without ideas. If they are not won over to a socialist world view by the intervention of conscious revolutionaries, they will continue to accept the bourgeois ideology of existing society. This is all the more likely because it is an ideology that flavours all aspects of life at present and is perpetuated by all media. Even were some workers ``spontaneously'' to come to a fully fledged scientific standpoint they would still have to argue with others who had not.

\blockquote{To forget the distinction between the vanguard and the whole of the masses gravitating towards it, to forget the vanguard's constant duty of raising ever-wider sections to its own advanced level, means simply to deceive oneself, to shut one's eyes to the immensity of our tasks, and to narrow down these tasks.\footnote{V. I. Lenin, \emph{Left-Wing Communism: An Infantile Disorder}.}}

This argument is not one that can be restricted to a particular historical period. It is not one, as some people would like to argue, that applies to the backward Russian working class of 1902 but not to those in the advanced nations today. The absolute possibilities for the growth of working-class consciousness may be higher in the latter, but the very nature of capitalist society continues to ensure a vast unevenness within the working class. To deny this is to confuse the revolutionary \emph{potential} of the working class with its present situation. As Lenin writes against the Mensheviks (and Rosa Luxemburg!) in 1905:

\blockquote{Use fewer platitudes about the development of the independent activity of the workers---the workers display no end of independent revolutionary activity which you do not notice!---but see to it rather that you do not demoralise undeveloped workers by your own tailism.\footnote{V. I. Lenin, \emph{Collected Works}, Vol.~9.}}

\blockquote{There are two sorts of independent activity. There is the independent activity of a proletariat that possesses revolutionary initiative, and there is the independent activity of a proletariat that is undeveloped and held in leading strings\ldots There are Social Democrats to this day who contemplate with reverence the second kind of activity, who believe they can evade a direct reply to pressing questions of the day by repeating the word ``class'' over and over again.\footnote{Ibid.}}

In short: stop talking about what the class as a whole can achieve, and start talking about how we as part of its development are going to act. As Gramsci writes:

\blockquote{Pure spontaneity does not exist in history: it would have to coincide with pure mechanical action. In the ``most spontaneous'' of movements the elements of ``conscious direction'' are only uncontrollable\ldots There exists a multiplicity of elements of conscious direction in these movements, but none of them is predominant\ldots\footnote{Gramsci, \emph{The Prison Notebooks}.}}

\blockquote{Man is never without some conception of the world. He never develops apart from some collectivity. For his conception of the world a man always belongs to some grouping, and precisely to that of all the social elements who share the same way of thinking and working.\footnote{Ibid.}}

Unless he is involved in a constant process of criticism of his world view so as to bring it coherence:

\blockquote{He belongs simultaneously to a multiplicity of men-masses, his own personality is made up in a queer way. It contains elements of the caveman and principles of the most modern advanced learning, shabby prejudices of all past historical phases, and intuitions of a future philosophy of the human race united all over the world.\footnote{Ibid.}}

\blockquote{The active man of the masses works practically, but does not have a clear theoretical consciousness of his actions, which is also a knowledge of the world insofar as he changes it. Rather his theoretical consciousness may be opposed to his actions. We can almost say that he has two theoretical consciousnesses (or one contradictory consciousness), one implicit in his actions, which unites him with all his colleagues in the practical transformation of reality, and one superficially explicit or verbal which he has inherited from the past and which he accepts without criticism\ldots (This division can reach the point) where the contradiction within his consciousness will not permit any action, any decision, any choice, and produces a state of moral and political passivity.\footnote{Ibid.}}

In this sense the question as to the preferability of ``spontaneity'' or ``conscious direction'' becomes that of whether it is:

\blockquote{preferable to think without having a critical awareness, in a disjointed and irregular way, in other words to ``participate'' in a conception of the world ``imposed'' mechanically by external environment, that is by one of the many social groups in which everyone is automatically involved from the time he enters the conscious world, or is it preferable to work out one's own conception of the world consciously and critically.\footnote{Ibid.}}

Parties exist in order to act in this situation---to propagate a particular world view and the practical activity corresponding to it. They attempt to unite together into a collectivity all those who share a particular world view and to spread this. They exist to give homogeneity to the mass of individuals influenced by a variety of ideologies and interests. But they can do this in two ways. The first Gramsci characterises as that of the Catholic Church. This attempts to bind a variety of social classes and strata to a single ideology by an iron discipline over the intellectuals that reduces them to the level of the ``ordinary people''. ``Marxism is antithetical to this Catholic position.'' Instead it attempts to unite intellectuals and workers so as to constantly raise the level of consciousness of the masses, so as to enable them to act truly independently. This is precisely why Marxists cannot merely ``worship'' the spontaneity of the masses: this would be to copy the Catholics in trying to impose on the most advanced sections the backwardness of the least.

For Gramsci and Lenin this means that the party is constantly trying to make its newest members rise to the level of understanding of its oldest. It has always to be able to react to the ``spontaneous'' developments of the class, to attract those elements that are developing a clear consciousness as a result of these.

\blockquote{To be a party of the masses not only in name, we must get ever-wider masses to share in all party affairs, steadily to elevate them from political indifference to protest and struggle, from a general spirit of protest to an adoption of Social-Democratic views, from adoption of these views to support of the movement, from support to organised membership in the party.\footnote{V. I. Lenin, \emph{Collected Works}, Vol.~10.}}

The party able to fulfil these tasks will not, however, be the party that is necessarily ``broadest''. It will be an organisation that combines, with a constant attempt to involve in its work ever wider circles of workers, a limitation on its membership to those willing to seriously and scientifically appraise their own activity and that of the party generally. This necessarily means that the definition of what constitutes a party member is important. The party is not to be made up of just anybody who wishes to identify himself as belonging to it, but only those willing to accept the discipline of its organisations. In normal times the numbers of these will be only a relatively small percentage of the working class; but in periods of upsurge they will grow immeasurably.

There is an important contrast here with the practice in Social-Democratic parties. Lenin himself realises this only insofar as Russia is concerned prior to 1914, but his position is clear. He contrasts his aim---a ``really iron strong organisation'', a ``small but strong party'' of ``all those who are out to fight''---with the ``sprawling monster, the new \emph{Iskra} motley elements of the Mensheviks''.\footnote{V. I. Lenin, \emph{Collected Works}, Vol.~7.} This explains his insistence on making a principle out of the question of the conditions for membership of the party when the split with the Mensheviks occurred.

Within Lenin's conception those elements that he himself is careful to regard as historically limited and those of general application must be distinguished. The former concern the stress on closed conspiratorial organisations and the need for careful direction from the top down of party officials, etc. ``Under conditions of political freedom our party will be built entirely on the elective principle. Under the autocracy this is impracticable for the collective thousands of workers that make up the party.''\footnote{V. I. Lenin, \emph{Collected Works}, Vol.~7.}

Of much more general application is the stress on the need to limit the party to those who are going to accept its discipline. It is important to stress that for Lenin (as opposed to many of his would-be followers) this is not a blind acceptance of authoritarianism. The revolutionary party exists so as to make it possible for the most conscious and militant workers and intellectuals to engage in scientific discussion as a prelude to concerted and cohesive action. This is not possible without general participation in party activities. This requires clarity and precision in argument combined with organisational decisiveness. The alternative is the ``marsh''---where elements motivated by scientific precision are so mixed up with those who are irremediably confused as to prevent any decisive action, effectively allowing the most backward to lead. The discipline necessary for such a debate is the discipline of those who have ``combined by a freely adopted decision''.\footnote{V. I. Lenin, \emph{One Step Forward, Two Steps Back}.} Unless the party has clear boundaries and unless it is coherent enough to implement decisions, discussion over its decisions, far from being ``free'', is pointless. Centralism for Lenin is far from being the opposite of developing the initiative and independence of party members; it is the precondition of this. It is worth noting how Lenin summed up the reasons for his battle for centralism over the previous two years in 1905. Talking of the role of the central organisation and of the central paper he says that the result was to be the:

\blockquote{creation of a network of agents\ldots that\ldots would not have to sit round waiting for the call to insurrection, but would carry out such regular activity that would guarantee the highest probability of success in the event of an insurrection. Such activity would strengthen our connections with the broadest masses of the workers and with all strata that are discontented with the aristocracy\ldots Precisely such activity would serve to cultivate the ability to estimate correctly the general political situation and, consequently, the ability to select the proper moment for the uprising. Precisely such activity would train all local organisations to respond simultaneously to the same political questions, incidents, and events that agitate the whole of Russia and to react to these ``incidents'' in the most rigorous, uniform and expedient manner possible\ldots\footnote{V. I. Lenin, \emph{Collected Works}, Vol.~9.}}

By being part of such an organisation worker and intellectual alike are trained to assess their own concrete situation in accordance with the scientific socialist activity of thousands of others. ``Discipline'' means acceptance of the need to relate individual experience to the total theory and practice of the party. As such it is not opposed to, but a necessary prerequisite of the ability to make independent evaluations of concrete situations. That is also why ``discipline'' for Lenin does not mean hiding differences that exist within the party, but rather exposing them to the full light of day so as to argue them out. Only in this way can the mass of members make scientific evaluations. The party organ must be open to the opinions of those it considers inconsistent.

\blockquote{It is necessary in our view to do the utmost---even if it involves certain departures from tidy patterns of centralism and from absolute obedience to discipline---to enable these grouplets to speak out and give the whole Party the opportunity to weigh the importance or unimportance of those differences and to determine where, how and on whose part inconsistency is shown.\footnote{V. I. Lenin, \emph{Collected Works}, Vol.~16.}}

In short, what matters is that there is political clarity and hardness in the party so as to ensure that all its members are brought into its debate and understand the relevance of their own activity. That is why it is absurd, as the Mensheviks tried to do, and as some people still do, to confuse the party with the class. The class as a whole is constantly engaged in unconscious opposition to capitalism; the party is that section of it that is already conscious and unites to try to give conscious direction to the struggle of the rest. Its discipline is not something imposed from the top downwards, but rather something that is voluntarily accepted by all those who participate in its decisions and act to implement these.

\paragraph{The Social-Democratic Party, the Bolshevik Party and the Stalinist Party}

We can now see the difference between the party as Lenin conceived it and the Social-Democratic party simultaneously envisaged and feared by Rosa Luxemburg and Trotsky. The latter was thought of as a party of the whole class. The coming to power of the class was to be the party taking power. All the tendencies within the class had to be represented within it. Any split within it was to be conceived of as a split within the class. Centralisation, although recognised as necessary, was feared as a centralisation over and against the spontaneous activity of the class. Yet it was precisely in this sort of party that the ``autocratic'' tendencies warned against by Luxemburg were to develop most. For within it the confusion of member and sympathiser, the massive apparatus needed to hold together a mass of only half politicised members in a series of social activities, led to a toning down of political debate, a lack of political seriousness, which in turn reduced the ability of the members to make independent political evaluations, increased the need for apparatus-induced involvement. Without an organisational centralisation aimed at giving clarity and decisiveness to political differences, the independence of the rank-and-file members was bound to be permanently undermined. Ties of personal affection or of deference to established leaders become more important than scientific, political evaluation. In the marsh, where no one takes a clear road (even if the wrong one), then there is no argument as to which is the right one. Refusal to relate organisational ties to political evaluations, even if done under the noble intention of maintaining a ``mass party'', necessarily led to organisational loyalties replacing political ones. This in turn entailed a failure to act independently given opposition from old colleagues (the clearest example of this tendency was undoubtedly Martov in 1917).

It is essential to understand that the Stalinist party is not a variant of the Bolshevik party. It too was dominated by organisational structures. Adherence to the \emph{organisation} rather than to the politics of the organisation mattered. Theory existed to justify an externally determined practice, not vice-versa. Organisational loyalties of the apparatus are responsible for political decisions (the former relate in turn to the needs of the Russian state apparatus). It is worth noting that in Russia a real victory of the apparatus over the party required precisely the bringing into the party of hundreds of thousands of ``sympathisers'', a dilution of the ``party'' by the ``class''. At best politically unsure of themselves, the ``Lenin levy'' could be relied upon to defer to the apparatus. The Leninist party does not suffer from this tendency to bureaucratic control precisely because it restricts its membership to those willing to be serious and disciplined enough to take \emph{political} and \emph{theoretical} issues as their starting point, and to subordinate all their activities to these.

But does this not imply a very elitist conception of the party? In a sense it does, although this is not the fault of the party, but of life itself, which gives rise to an uneven development of working-class consciousness. The party to be effective has to aim at recruiting all those it conceives of as being most ``advanced''. It cannot reduce its own level of science and consciousness merely in order not to be an ``elite''. It cannot, for instance, accept that chauvinist workers are ``as good as'' internationalist party members, so as to take account of the ``self-activity'' of the class. But to be a ``vanguard'' is not the same as to substitute one's own desires, or policies, or interests for those of the class.

For Lenin the party is not the embryo of the workers' state---the workers' council is. The working class as a whole will be involved in the organisations that constitute its state, the most backward as well as the most progressive elements. ``Every cook will govern.'' In Lenin's major work on the state, the party is hardly mentioned. The function of the party is not to be the state, but rather to carry out continual agitation and propaganda among more backward elements of the class so as to raise their self consciousness and self reliance to the pitch that they will both set up workers' councils and fight to overthrow the forms of organisation of the bourgeois state. The Soviet state is the highest concrete embodiment of the self-activity of the whole working class; the party is that section of the class that is most conscious of the world historical implications of this self-activity.

The functions of the workers' state and of the party should be quite different (which is why there can be more than one party in a workers' state). One has to represent all the diverse interests of all the sections---geographical, industrial, etc.---of the workers. It has to recognise in its mode of organisation all the heterogeneity of the class. The party, on the other hand, is built around those things that unite the class nationally and internationally. It constantly aims, by ideological persuasion, to overcome the heterogeneity of the class. It is concerned with national and international political principles, not parochial concerns of individual groups of workers. It can only persuade, not coerce, these into accepting its lead. An organisation that is concerned with participating in the revolutionary overthrow of capitalism by the working class cannot conceive of substituting itself for the organs of direct rule of that class. Such a perspective is only available to the Social-Democratic or Stalinist Party (and both have been too afraid of mass self-activity to attempt this substitution through revolutionary practice in advanced capitalist countries). Existing under capitalism, the revolutionary organisation will of necessity have a quite different structure to that of the workers' state that will arise in the process of overthrowing capitalism. The revolutionary party will have to struggle within the institution of the workers' state for its principles as against those with opposed ones; this is only possible because it itself is not the workers' state.

Lenin's theory of the party and his theory of the state are not two separate entities, capable of being dealt with in isolation from one another. Until he developed the theory of the state, he tended to regard the Bolshevik Party as a peculiar adaptation to Russian circumstances. Given the Social-Democratic (and later the Stalinist) conception of the party becoming the state, it is only natural for genuinely revolutionary and therefore democratic socialists not to want to restrict the party to the most advanced sections of the class, even if the need for such an organisation of the most conscious sections is recognised. This explains Rosa Luxemburg's ambiguity over the question of political organisation and theoretical clarity. It enables her to counterpose the ``errors committed by a truly revolutionary movement'' to the ``infallibility of the cleverest central committee''. But if the party and the institutions of class power are distinct (although one attempts to influence the other), the ``infallibility'' of the one is a central component in the process by which the other learns from its errors. It is Lenin who sees this. It is Lenin who draws the lessons, not (at least until the very end of her life) Luxemburg. The need is still to build an organisation of revolutionary Marxists that will subject their situation and that of the class as a whole to scientific scrutiny, will ruthlessly criticise their own mistakes, and will, while engaging in the everyday struggles of the mass of workers, attempt to increase their independent self-activity by unremittingly opposing their ideological and practical subservience to the old society.

A reaction against the identification of class and party elite made by both Social Democracy and Stalinism is very healthy. It should not, however, prevent a clear-sighted perspective of what we have to do to overcome their legacy.
